\chapter{画像と動画の認識}
\chaptermark{認識}
\label{sec:recognition}

\section{課題：画素が違っても同じもの}

第\ref{sec:sensors}章では、画像をマシンの入力まで運んだ。目の約100万個の出力ニューロンと、主にエッジと変化を伝える圧縮された流れである。だがセンサと行動の間には、これまで触れずにきた段階がある。絵の中のネコは画素の集まりではない。半フレームずらし、回転させ、遠ざけ、照明を半分消せば、ほとんどすべての画素が変わる。それでも答えは「ネコ」のままでなければならない。エンジニアはこれを\emph{不変性}と呼ぶ。分類器の答えは、平行移動、拡大縮小、回転、照明によって変わってはならない。

難しさは簡単な実験でよく分かる。$24$ 点からなる1次元の「網膜」と、どこにでも置ける5点ずつの二つの図形を考える。1か所で記憶したテンプレートと入力を比べる分類器の正答率は $49$--$52\%$ で、コイン投げと変わらない。ずれがあると、画素ごとの比較は完全に破綻する。別の構成が必要である。

\section{段のパイプライン}

脳がこれをどれほど速く行うかは、第\ref{sec:code}章ですでに知っている。一瞬だけ示された絵の中の動物はおよそ $150\ms$ で認識され、信号は十数段を1段あたり $10$--$20\ms$ で通過する~\cite{Thorpe1996}。各段でニューロンが出せるのはスパイク0個か1個である。つまり速い認識とは、長い熟考も繰り返しもない、パイプラインの1回の通過である。問題は各段が何をするかである。

視覚の最初の段での測定~\cite{HubelWiesel1962}が示した答えは、交互に現れる二つの演算からなる（図~\ref{fig:conveyor}）。

\emph{選択性。} $S$ 段のニューロンは入力の小さな領域を見て、そこに特定の部品の組み合わせ（このエッジと、その隣のあのエッジ）があるときだけ発火する。これは部品についての AND であり、どの部品一つが与える入力よりも閾値が高い、第\ref{sec:glutcomputer}章のしきいニューロンである。

\emph{不変性。} $C$ 段のニューロンは、同じ組み合わせを異なる場所で探す多数の $S$ ニューロンの出力を集め、それがどこか一か所でも見つかれば発火する。これは位置についての OR、より正確には最大値の選択である。

1次元モデルでは、二つの演算は次のように書ける：
\begin{empheq}[box=\keybox]{equation}
s_{f,p} \;=\; \Bigl[\sum_i t_{f,i}\,x_{p+i} - \theta\Bigr]_+ ,
\qquad
c_f \;=\; \max_p\, s_{f,p} ,
\label{eq:scstage}
\end{empheq}
ここで $x$ は入力、$t_{f,i}$ は特徴 $f$ のテンプレート、$p$ は位置、$\theta$ は閾値である。第1式は応答を選択的にし、第2式は位置に依存しないようにする。多数の入力の最大値は、第\ref{sec:gaba}章の道具で得られる。相互抑制は最も強い入力を残し（命題~\ref{prop:wta}）、全体の活動による除算は明るさの違いによる応答の差をならす~\cite{Carandini2012}。

\begin{figure}[t]
\centering
\begin{tikzpicture}[>=Stealth,
  px/.style={draw,minimum size=0.36cm,inner sep=0pt},
  on/.style={px,fill=blue!55!black},
  snode/.style={circle,draw,very thick,minimum size=0.62cm,inner sep=0pt,font=\scriptsize,fill=blue!6},
  cnode/.style={circle,draw,very thick,minimum size=0.8cm,inner sep=0pt,font=\scriptsize,fill=red!10}]
% two inputs: the same shape at two positions
\foreach \row/\sh/\lab in {0/1/{フレーム1},1/7/{フレーム2}}{
  \begin{scope}[yshift=-\row*1.0cm]
  \node[left,font=\scriptsize] at (-0.25,0) {\lab};
  \foreach \i in {0,...,13}{\node[px] at (\i*0.4,0) {};}
  \foreach \d in {0,1,3,4}{\node[on] at ({(\sh+\d)*0.4},0) {};}
  \end{scope}}
% S stage: detectors of the shape at several positions
\foreach \k in {0,...,5}{
  \node[snode] (S\k) at ({0.8+0.8*\k},1.7) {$S$};}
\node[font=\scriptsize,left] at (-0.25,1.7) {$S$ 段：部品の AND};
\draw[->,thick,blue!60!black] (1.2,0.25) -- (S1.south);
\draw[->,thick,orange!85!black] (3.6,0.25) -- (S4.south);
\node[font=\scriptsize,blue!60!black,anchor=east] at (1.25,0.85) {フレーム1};
\node[font=\scriptsize,orange!85!black,anchor=west] at (3.9,0.85) {フレーム2};
% C stage
\node[cnode] (C) at (2.8,3.25) {$C$};
\node[font=\scriptsize,left] at (-0.25,3.25) {$C$ 段：位置の OR};
\foreach \k in {0,...,5}{\draw[->,gray!60!black] (S\k.north) -- (C);}
\draw[->,very thick,red!70!black] (C.north) -- ++(0,0.6) node[above,font=\scriptsize,black] {「図形あり」：どちらのフレームでも};
% next stages
\draw[decorate,decoration={brace,amplitude=5pt},gray!60!black] (6.3,1.4) -- (6.3,-1.3);
\node[right,font=\scriptsize,align=left] at (6.5,0.05) {同じ物体、\\違う画素};
\draw[->,thick,densely dashed,gray!60!black] (3.6,3.25) -- (5.9,3.25) node[right,font=\scriptsize,black,align=left] {次の $S$, $C$ の対：\\より大きく、より複雑、\\より不変};
\end{tikzpicture}
\caption{パイプラインの一対の段。$S$ ニューロンは同じ部品の組み合わせを異なる場所で探し、$C$ ニューロンはそれがどこかで見つかれば応答する~\eqref{eq:scstage}。フレーム1とフレーム2の図形は異なる $S$ ニューロンを興奮させるが、同じ $C$ を興奮させる。次の段の対は、同じことをより大きく複雑な組み合わせについて繰り返す。}
\label{fig:conveyor}
\end{figure}

同じ1次元モデルで、一対の段~\eqref{eq:scstage}は図形をどの位置でも誤りなく認識する。入力の各点を確率 $0.1$ でランダムに反転させると正答率は $86\%$、確率 $0.2$ では $70\%$ である。コイン投げは相変わらず半分である。このような対をいくつか積み重ねてみよう。最初の対はエッジを、次の対はエッジの組み合わせを、さらに上は物体の部分を、最上段は物体全体を探す。対を重ねるごとに組み合わせは大きくなり、答えは物体がどこにどう置かれているかにますます依存しなくなる~\cite{Riesenhuber1999}。

\begin{keyblock}
\begin{proposition}[選択性と不変性のパイプライン]
\label{prop:conveyor}
速い認識は十数段の1回の通過である。各段の対で、部品の AND~\eqref{eq:scstage}が応答を選択的にし、位置の OR がずれに依存しないようにする。この二つの演算を交互に繰り返すと、物体を区別しつつ、それがどこにどう見えるかにほとんど依存しない答えが得られる。最初の段の仕組みと通過の速さは測定されている。連鎖全体がこの二つの演算に帰着するというのは単純化である。
\end{proposition}
\end{keyblock}

この仕組みに見覚えがあるなら、それで正しい。まさにこの交互の繰り返し、つまりすべての場所で特徴を探し、次に近くの場所どうしをまとめることを、エンジニアは人工の\emph{畳み込み}ネットワークに移した~\cite{Fukushima1980}。最近は逆向きの流れもある。物体認識を学習した畳み込みネットワークが、視覚の上位段のニューロン応答について知られている最良のモデルであることが分かった~\cite{Yamins2014}。最上段の結果は簡単に読み出せる。最終段の約100個のニューロンの発火率の合計から、線形の判定器が提示後0.1秒で物体を認識する~\cite{Hung2005}。これは第\ref{sec:code}章の集団のレート符号である。

\section{いない教師：動画から学ぶ}

畳み込みネットワークは、ラベルの付いた数百万枚の画像で学習する。脳にラベルを与えてくれる者はほとんどいない。子どもは何時間も物を見るが、「カップ」という言葉を聞くのはたまにである。では $C$ 段は、どの $S$ ニューロンをまとめるべきかをどうやって知るのか。「ここにあるこの図形」と「あそこにあるこの図形」をまとめるには、それが同じ図形だとすでに知っていなければならないはずである。

ヒントは動画そのものにある。世界は連続的に変わる。視野の中の物体は、ずれ、回転し、照明が変わる間も同じ物体であり続け、視線が別の物体に移るのはときどきにすぎない。\emph{ゆっくり}変わるものはおそらく物体に属し、速く変わるものはその位置と見え方に属する~\cite{WiskottSejnowski2002}。したがって $C$ ニューロンには、第\ref{sec:dofa}章の適格度トレース付きヘッブ則で十分である。同時に来たものだけでなく、続けて来たものを結びつければよい~\cite{Foldiak1991}：
\begin{empheq}[box=\keybox]{equation}
\tau_{\text{tr}}\,\frac{\dd\bar y_k}{\dd t} \;=\; -\bar y_k + y_k ,
\qquad
\frac{\dd\mathbf w_k}{\dd t} \;=\; \eta\,\bar y_k\,\bigl(\mathbf x - \mathbf w_k\bigr) .
\label{eq:traceinvariance}
\end{empheq}
ここで $y_k$ は抑制による競合に勝った $C$ ニューロン $k$ の活動、$\bar y_k$ はそのトレースで、規則~\eqref{eq:threefactor}と同じRC回路である。$\mathbf x$ は $S$ ニューロンの出力である。物体の最初の見え方で勝ったニューロンは、2番目の見え方が来たときにもまだそれを覚えていて、それも自分に引き寄せる。

\begin{figure}[t]
\centering
\begin{tikzpicture}[>=Stealth,x=1.0cm,y=1.0cm]
\foreach \i/\lab in {0/1,1/2,2/3,3/4}{
  \node[draw,thick,fill=blue!8,minimum width=0.9cm,minimum height=0.55cm,font=\scriptsize] at (\i+0.5,2.2) {$A$@\lab};}
\foreach \i/\lab in {0/4,1/3,2/2,3/1}{
  \node[draw,thick,fill=orange!12,minimum width=0.9cm,minimum height=0.55cm,font=\scriptsize] at (\i+6.5,2.2) {$B$@\lab};}
\node[font=\scriptsize,gray!40!black] at (5.0,2.2) {休止};
\draw[->,gray!70!black] (0,0) -- (10.4,0) node[below left,font=\scriptsize,black] {時間};
% traces
\draw[thick,blue!60!black] (0,0.05) .. controls (0.4,1.2) and (0.8,1.3) .. (1.2,1.35) -- (4,1.4) .. controls (4.6,0.5) and (5.2,0.15) .. (6,0.08) -- (10,0.05);
\draw[thick,orange!85!black] (0,0.05) -- (6,0.05) .. controls (6.4,1.2) and (6.8,1.3) .. (7.2,1.35) -- (10,1.4);
\node[font=\scriptsize,blue!60!black,anchor=west] at (1.4,1.65) {トレース $\bar y_1$：$A$ の通過中ずっと保たれる};
\node[font=\scriptsize,orange!85!black,anchor=west] at (7.2,1.65) {トレース $\bar y_2$};
\draw[decorate,decoration={brace,amplitude=4pt},gray!60!black] (4.0,2.6) -- (0,2.6);
\draw[decorate,decoration={brace,amplitude=4pt,mirror},gray!60!black] (0,-0.35) -- (4,-0.35) node[midway,below=4pt,font=\scriptsize,black] {$A$ のすべての見え方がニューロン1に結びつく};
\end{tikzpicture}
\caption{動画が不変性をどう教えるか（模式図）。物体 $A$ が四つの位置を通過し（$A$@1 は位置1の $A$）、休止をはさんで物体 $B$ が続く。$A$ の最初の見え方で勝ったニューロンのトレース~\eqref{eq:traceinvariance}は通過の間ずっと保たれ、$A$ のすべての見え方が一つのニューロンに結びつく。休止がトレースを消し、$B$ は別のニューロンに割り当てられる。}
\label{fig:tracevideo}
\end{figure}

これをモデルで確かめた（図~\ref{fig:tracevideo}）。二つの $C$ ニューロンが $16$ 個の $S$ ニューロン（二つの図形×八つの位置）からの入力を奪い合う。図形はすべての位置を1位置あたり $50\ms$ で通過し、次に $0.3\,\text{s}$ の休止、そして次のランダムな図形が来る。ラベルは一切ない。トレースが短い $\tau_{\text{tr}}=5\ms$ のとき、ニューロンは入力を\emph{位置}で分け合い、答えが図形と一致する割合はコイン投げと変わらない。10回の試行の平均で $52\%$ である。トレースが一つの位置の提示時間と同程度、$\tau_{\text{tr}}\approx20\ms$ 以上になると、各ニューロンは自分の図形の\emph{すべての}位置を集める（$20$, $50$, $200\ms$ では10回の試行すべてで）。不変性が動画だけから学習されたのである。物体間の休止は重要である。休止がないと、トレースが次の物体に流れ込み、二つを混同させる。

同じことは実験でも見られる。眼があるところから別のところへ跳ぶ間に物体をこっそり差し替えると、上位段のニューロンは1〜2時間のうちに、差し替えた物体に元の物体と同じように応答し始める。続けて来たものを結びつけるのである~\cite{LiDiCarlo2008}。脳の中の不変性は、確かに時間の連続性から学習される。視覚の学習全体がこの規則でどこまで説明できるかは、仮説である。

\section{動画とは運動である}

動画は学習のためのフレームの流れであるだけでなく、それ自体が課題でもある。何が、どちらへ、どれだけ速く動いているか。最初の一歩はおなじみの、第\ref{sec:gaba}章の方向検出器である（図~\ref{fig:direction}）。そこでは遅れて来る抑制が一方向の運動を打ち消す。このような検出器が視野全体に並び、その先は形の特徴と同じ扱いを受ける。同じ方向の検出器を異なる場所から多数集める段は、大きな物体の運動に、それがどこにあっても応答する。これは同じ AND と OR の対~\eqref{eq:scstage}であり、特徴が「エッジ」ではなく「時間 $d$ の間にずれたエッジ」であるだけである。

動画はまた予測可能でもある。次のフレームは前のフレームとほとんど同じである。\emph{予測符号化}の仮説では、上位の段は下位の段に予測を送り、上へ送られるのは主に誤差、つまり予測が捉えきれなかった部分である~\cite{RaoBallard1999}。これは命題~\ref{prop:economy}と同じスパイクの節約である。すでに分かっていることではなく、ニュースに支払う。個々の観察はこの仮説と整合するが、パイプラインの一般的な仕組みとしては証明されていない。

\section{脳と畳み込みネットワーク}

類似はどこまで及ぶのか。一つの物体の速い認識については、確かに深い~\cite{Yamins2014}。だが脳には、単純な畳み込みネットワークにないものもある。

\emph{フィードバック。} 脳のパイプラインは順方向だけではない。どの段にも後ろ向きと横向きの結合がある。重なりや悪い照明のある難しい画像では、上位段の応答はより遅れて形成され、その遅い応答は順方向のネットワークよりフィードバックのあるネットワークのほうがよく記述する~\cite{Kar2019}。1回の通過は易しい場合を解き、難しい場合には何周かが必要になる。

\emph{頑健性。} 人工ネットワークは、目には見えない画素の改ざんでだませる~\cite{Szegedy2014}。人に見えない変化が「パンダ」を「テナガザル」に変える~\cite{Goodfellow2015}。ヒトの視覚はこのような改ざんにはるかに強いが、無敵ではない。多くのネットワークを同時にだます画像は、短時間提示ではヒトの選択もずらす~\cite{Elsayed2018}。頑健性がなぜこれほど違うのかは未解決の問題である。候補はノイズ、全体の活動による除算、フィードバックである。

\emph{教師。} ネットワークにはラベル付きの例が数百万必要だが、脳に要るのは主にラベルなしの動画~\eqref{eq:traceinvariance}と、何が重要かについての\nt{DOPA}からのわずかなヒントである。

\emph{エネルギー。} 脳全体で約 $20$ ワット（命題~\ref{prop:economy}）、認識はその一部にすぎない。学習済みの畳み込みネットワークはGPU上で数百ワットを使う。

\section{錯視：マシンはどこでなぜ間違えるか}
\label{sec:illusions}

他人の回路を理解したいエンジニアは、ふつうでない信号を入れて、どこで間違えるかを見る。視覚については、そうした信号はずっと前に考案されている。錯視である。錯視は故障ではない。どの錯視も、ふだんの状況では正しく働くマシンの特定の機構を指し示し、特別に選んだ画像でその機構が正体を現す。

\emph{妥当な期待。} 入力はノイズを含むので、マシンはそれを世界でふつう起こることで補う。遅い運動は速い運動より多い。入力が信頼できないとき、たとえば画像のコントラストが低いとき、速さの推定は遅いほうへずれ、淡い物体は明るい物体より遅く動いて見える~\cite{Weiss2002}。明るさの錯視の多くも同じように説明される。我々が見るのは画素の明るさではなく、その明るさが過去に最もよく対応していた表面である~\cite{Purves2011}。ある人には白と金に、別の人には青と黒に見えるドレスの写真も同じ機構である。画像はあいまいで、人は無意識に想定する照明の点で分かれる~\cite{LaferSousa2015,Wallisch2017}。個人差は測定されている。ここで脳が確率論の規則に従って入力と期待を組み合わせているというのは、十分な根拠のある仮説である。

\emph{ゲイン調整。} 滝を1分間見つめてから、動かない岩を見てみよう。岩が上へ流れていく。「下」と「上」のチャネルは\eqref{eq:dualrail}のような配線の対である。ゲイン調整~\eqref{eq:nakarushton}が疲れたチャネルを弱め、対のつり合いがずれたのである。周囲が中心の見え方を変える傾きやコントラストの錯視は、第\ref{sec:gaba}章と同じく、近傍の活動による除算で説明される~\cite{Schwartz2009}。慎重さを教える例もある。格子の白い帯の交差点に見える暗い斑点は、長く近傍の単純な抑制で説明されてきたが、格子を変形した実験によって、その説明が成り立たないことが示された~\cite{SchillerCarvey2005}。

\emph{遅延の補償。} 目から上位段までの経路には数十ミリ秒かかり、動く物体はその間に移動してしまう。その横で静止した点が光ると、両者は一致していたのに、物体は閃光より前にあるように見える~\cite{Nijhawan1994}。一つの説明では、マシンは遅延を補償するために運動を前に延長する。第\ref{sec:muscles}章と同じく、フィードバックが遅れるなら予測しなければならない。この錯視には説明がいくつかあり、論争は決着していない。

\emph{タイミング符号の誤り。} 色の急な移り変わりのある輪からなる静止画が、回転しているように見える。コントラストの高い部分は淡い部分より速く処理され、遅延の差~\eqref{eq:latency}を方向検出器が運動として読む~\cite{Conway2005}。錯視を起動するのは、目の小さな不随意の跳躍と瞬きである。跳躍のたびに入力が更新され、検出器は再び「運動」を見る~\cite{OteroMillan2012}。これは第\ref{sec:code}章のタイミング符号の読み違いである。

\emph{一つの中の二つの絵。} ある面をこちらに向けたり別の面を向けたりする立方体の枠や、左右の目に違う絵を見せた場合、知覚は数秒ごとに切り替わる。最良のモデルは、二つの解釈の相互抑制、遅い疲労、ノイズである~\cite{MorenoBote2007}。つまり第\ref{sec:muscles}章で歩行のリズムを生んだのと同じ仕組み~\eqref{eq:halfcentre}である。切り替わりのタイミングはノイズと疲労が決める。モデル~\cite{MorenoBote2007}では主役はノイズで、疲労は補助にすぎない。

では人工ネットワークはどうか。動画の次のフレームを予測するよう学習したネットワークは、同じ静止した輪に回転を見て、対照画像には見ない~\cite{Watanabe2018}。画像のノイズ除去と鮮鋭化を学習したネットワークは、ヒトと同じ色と明るさの錯視にかかる~\cite{GomezVilla2020}。つまり錯視の一部は、神経組織の特性ではなく、マシンが学習する課題そのものから生じる。どの錯視が脳だけに固有なのかは、どんな視覚モデルにとってもよい試験になる。

\begin{keyblock}
\begin{proposition}[機構の痕跡としての錯視]
\label{prop:illusions}
錯視は、ふだんの世界で役に立つ機構がふつうでない入力を受けたときに生じる。期待が信頼できない入力に勝ち、ゲイン調整がチャネルの対をずらし、遅延の補償が物体を前へ押し出し、遅延の差が運動として読まれ、ノイズと疲労をともなう相互抑制が切り替わる。どの錯視も、自分の機構を指し示す試験信号である。錯視そのものは測定されている。その説明は、十分な根拠のあるものから論争中のものまである。
\end{proposition}
\end{keyblock}

\section{まとめ}

\begin{table}[t]
\centering
\footnotesize
\setlength{\tabcolsep}{4pt}
\renewcommand{\arraystretch}{1.15}
\begin{tabular}{>{\raggedright}p{3.0cm}>{\raggedright}p{4.6cm}>{\raggedright\arraybackslash}p{5.2cm}}
\toprule
マシンは何をするか & どうやって & 分かっていること \\
\midrule
$150\ms$ で認識する & 十数段の1回の通過 & 速さは測定済み \\
応答を選択的にする & 部品の AND、$S$ 段~\eqref{eq:scstage} & 最初の段で測定済み \\
応答を不変にする & 位置の OR、$C$ 段；抑制と除算 & 最初の段で測定済み；上位段については単純化 \\
答えを出す & 最終段の約100個のニューロンの発火率、線形の読み出し & 測定済み \\
ラベルなしで学ぶ & 連続した動画上のトレース付きヘッブ則~\eqref{eq:traceinvariance} & 物体の差し替えが応答を変える（測定済み）；主要な規則としては仮説 \\
運動を見る & 方向検出器、次に同じ AND／OR の対 & 測定済み \\
予測できるものを節約する & 上へは予測誤差が送られる & 仮説 \\
難しい場合を解く & フィードバック、何周か & 遅い応答とフィードバックの役割は測定済み \\
予測どおりに間違える & 錯視：期待、ゲイン調整、遅延、ノイズと疲労をともなう相互抑制 & 錯視は測定済み；説明は根拠のあるものから論争中のものまで \\
\bottomrule
\end{tabular}
\caption{マシンは画像と動画をどう認識するか。}
\label{tab:recognition}
\end{table}

認識は、すでに手元にある部品で組み立てられている（表~\ref{tab:recognition}）。閾値が部品の AND を、相互抑制が位置の OR を、除算が明るさへの非依存を、トレース付きヘッブ則がいない教師の代わりを与える。エンジニアは視覚から選択性と不変性の交互の繰り返しを借り、それをもとに畳み込みネットワークを作った。逆方向には、脳はまだ肝心なもの、つまりラベルなしの動画からほとんどエネルギーを使わずに学ぶ能力を明け渡していない。

\section{問題}

\begin{exercise}[\lvlA]
\label{ex:12:1}
\emph{1段での発火率。} パイプラインの1段は $15\ms$ 続き、ニューロンは毎秒 $20$ スパイクで発火し、ノイズは $c=0.1$~\eqref{eq:correlated}で相関している。100個のニューロンの集団の発火率の相対誤差と、集団をいくら大きくしたときのその極限はいくらか。1段で区別できる発火率のレベルはいくつか。
\end{exercise}

\begin{exercise}[\lvlA]
\label{ex:12:2}
\emph{1回の認識のエネルギー。} $150\ms$ の1回の通過で、$10^7$ 個のニューロンからなる10段がそれぞれ高々1スパイクを出し、1スパイクのコストは $10^{-10}$~J である。(a)~この $150\ms$ の間の脳全体のエネルギーのうち、認識はどれだけの割合を占めるか。(b)~画像を $10\ms$ で認識する $300$~W のアクセラレータは何倍多く使うか。
\end{exercise}

\begin{exercise}[\lvlB]
\label{ex:12:3}
\emph{$S$ 段の閾値。} $24$ 点の「網膜」、図形 $A=11011$ と $B=10101$、テンプレート~\eqref{eq:scstage}は図形の1の位置で $+1$、0の位置で $-1$ とする。(a)~$S$ ニューロンが1点の反転を許容しつつ、他方の図形には沈黙する閾値の範囲はどこか。(b)~$100\times100$ の画像上で $5\times5$ の特徴10個に要る $S$ ニューロンの数はいくつか。
\end{exercise}

\begin{exercise}[\lvlB]
\label{ex:12:4}
\emph{二つの絵の一方はどれだけ続くか。} 切り替えの仕組み~\eqref{eq:halfcentre}で $\tau\ll\tau_a$ とし、グループ1が勝っている間は $r_2=0$, $r_1=I-a_1$ とする。(a)~$I=1$, $\beta=2$, $b=2$, $\tau_a=0.4\,\text{s}$ で勝利の持続時間を求めよ。(b)~絵が3秒ごとに切り替わるのはどの $\tau_a$ か。
\end{exercise}

\begin{exercise}[\lvlB]
\label{ex:12:5}
\emph{シミュレーション：不変性の代償。} \texttt{models/\allowbreak recognition\_models.py} の関数 \texttt{two\_stage}（$4000$ 試行）を使い、$N=24$ で段の対 $S$, $C$ が4回に1回間違える点の反転確率 $p$ を求めよ。$p=0.1$ で、$N=8$ から $N=192$ まで正答率はどう変わるか。
\end{exercise}

\begin{exercise}[\lvlB]
\label{ex:12:6}
\emph{シミュレーション：トレースの窓。} \texttt{models/\allowbreak recognition\_models.py} の関数 \texttt{trace\_learning} を、休止を $0.3\,\text{s}$ として10回実行し、$5\ms$ から $2\,\text{s}$ までのどの $\tau_{\text{tr}}$ ですべての実行が誤りなく不変性を学習するかを求めよ。この窓の上限はどこから来るか。
\end{exercise}

\begin{exercise}[\lvlB]
\label{ex:12:7}
\emph{どこででも運動を。} 間隔 $0.1^\circ$ の点の列の隣り合う点に、第\ref{sec:gaba}章の方向検出器を置く。遅延 $d$ の $A$ からの抑制は、$B$ からの興奮とのずれが $5\ms$ 以内なら興奮を打ち消す。その上に $C$ 段がある。毎秒 $5$ から $20^\circ$ の逆向きの運動で $C$ を沈黙させるには、どの遅延が必要か。
\end{exercise}

\begin{exercise}[\lvlC]
\label{ex:12:8}
\emph{画素の改ざん。} きれいな図形について \texttt{two\_stage} モデル（$N=24$）の答えを変えるには、何点反転すればよいか。同じくずれに不変な分類器「入力の1が $3.5$ 個より多ければ $A$」ではどうか。$p=0.1$ でのその正答率は、段の対の $0.86$ と比べてどうか。
\end{exercise}
